% Page 021

\seite{21}

\jahresueberschrift{Das Jahr 1893}
\pstart
Die Schülerzahl betrug in diesem Jahre 54.\ob
Die Gemeinde hatte es übernommen, den Communikationsweg\ob
zwischen Sefferweich und Seffern zum großen Teile\ob
abzutragen und tiefer zu legen. Die Arbeit ist im Herbst\ob
dieses Jahres in Angriff genommen worden. Auch hier war\ob
ein großer Geldzuschuß staatlicherseits in Aussicht gestellt.\ob
Die Arbeit war dem Taglöhner Fried. Pardon für 1000 M.\ob
übertragen worden. Nachdem dieser ca 800 M.\ der\ob
Summe bezogen hatte und kaum ⅓ der Arbeit vollendet\ob
hatte, ließ er die Arbeit im Stich. Leider hatte er\ob
keinen solventen Bürgen, sodaß man beiden Beteiligten\ob
nichts anhaben konnte. Die Gemeinde mußte nun auf\ob
die Klage des Fortsprungs zu Werke gehen. Die versprochenen\ob
Gelder sind bis jetzt noch nicht eingetroffen. Für die Land\ohb
entschädigung muß die Gemeinde aufkommen.

In der Pfarrei wurde in diesem Jahre das hl. Sakrament der\ob
Firmung gespendet.

Im Herbste wurde in dem Gelände von Sefferweich das\ob
Manöver abgehalten. Die Landentschädigungen sind zur Zufrieden\ohb
heit der Bürger ausgefallen. Mit diesem Jahre hörte der\ob
hohe Preis des schwedischen Klees auf. Während man früher\ob
1--1,50 M.\ pro ½ Kilo erhielt, ist der Preis auf 25--30 Pfg.\ob
gesunken. Wegen Futtermangel war teilweise halbtags\ob
Schule eingeführt worden. Durch die Gemeinde wurde einige\ob
hundert Centner Kraftfutter bezogen, die Bezahlung ward\ob
auf drei Jahre verteilt.
\pend
